% Cover letter for the ASEJ submission.
% cover_letter_ASEJ.md is the editable source of record; this file mirrors it so a
% PDF exists for portals that want the letter as an attachment rather than pasted
% into a text box. Keep the two in step.
\documentclass[11pt,a4paper]{article}
\usepackage[T1]{fontenc}
\usepackage[utf8]{inputenc}
\usepackage{mathptmx}
\usepackage[margin=2.4cm]{geometry}
\usepackage{microtype}
\usepackage{array}
\pagestyle{empty}
\setlength{\parindent}{0pt}
\setlength{\parskip}{7pt}

\begin{document}

To the Editor-in-Chief\\
Ain Shams Engineering Journal\\
editor@eng.asu.edu.eg

Dear Editor,

We submit for your consideration the original research article \textbf{``Predictive fidelity versus
training utility in surrogate-based reinforcement learning for HVAC control: a role-separated hybrid
remedy''} for publication in \emph{Ain Shams Engineering Journal}.

\textbf{What the paper does.} Heating, ventilation and air-conditioning systems are the largest
single energy consumer in buildings, and reinforcement-learning (RL) supervisory controllers are
almost always trained against a fast learned surrogate of the building rather than against the
building itself. The whole practice rests on an assumption that, to our knowledge, has never been
tested directly: that a surrogate with higher predictive fidelity is a better training environment.
Using the open BOPTEST benchmark we show it is false, diagnose why, and provide an engineering
remedy.

\textbf{Main results.}

\begin{enumerate}\itemsep2pt \parskip0pt
\item A calibrated grey-box resistance-capacitance twin is by far the better predictor (24 h
rollout error $0.644\,^{\circ}$C against $1.557\,^{\circ}$C for a compact black-box surrogate), yet
the controller trained directly on it collapses on the live emulator on every seed, with comfort
violation above 77\,\%.
\item A matched-resolution ablation, in which the \emph{same} black-box architecture is retrained on
a finer corpus, becomes more accurate ($0.876\,^{\circ}$C) and equally unusable. This closes the
obvious model-class confound and identifies the coarse per-step temperature increment, that is the
temporal resolution of the training corpus, as the operative variable.
\item A role-separated hybrid training environment, in which the black-box surrogate supplies the
rollout dynamics and a frozen physical twin acts only as a per-step model-disagreement reward
censor, gives the best cross-window robustness of the study (maintenance score 0.041, comfort
violation below 5\,\%) at about 85 times live-simulator throughput.
\item Two boundary tests bound the recipe: the censor weight is controller-family specific, and
transfer to three hydronic test cases is component-level, since the calibration pipeline transfers
while the frozen policy does not.
\end{enumerate}

\textbf{Fit to the journal.} The work sits at the intersection of mechanical and control engineering
with computational methods, which the journal explicitly welcomes (``papers in which knowledge from
other disciplines is integrated with engineering''). It is fully reproducible on an open,
standardized building-control benchmark, and its practical recommendation is directly actionable for
building-control practitioners: a surrogate intended as an RL training environment should be
validated on downstream closed-loop utility and on the smoothness of its action-response surface,
not on predictive error alone.

\textbf{Compliance.} The manuscript is prepared for double-anonymized review: a separate title page
carries all author, affiliation, funding and interest information, and the manuscript file itself
contains no identifying material. All figures and tables are placed next to the relevant text with
captions below each object, sections are numbered, references are numbered in order of appearance
with journal titles abbreviated, and six keywords are given. Supplementary Material with the full
reproduction detail is attached separately.

\textbf{A note on the manuscript format.} We have submitted in Elsevier's two-column layout rather
than the 12\,pt single-spaced single-column format named in the Guide for Authors, and we would like
to explain the choice rather than leave it to inference. Read literally, the two instructions cannot
be satisfied together: a recent Full Length Article in this journal runs to roughly 10\,000 words
with 9 figures and 7 tables, which in 12\,pt single-column would occupy some 26 pages, well beyond
the stated 10-page limit. We therefore took the page limit to refer to the journal's own two-column
layout, in which our manuscript occupies 12 pages with 9 figures, 5 tables and 55 references, close
to the length of a typical article in the journal. We would of course be glad to supply the
manuscript in the single-column 12\,pt format instead if the editorial office prefers it.

The work is original, has not been published previously and is not under consideration elsewhere.
All authors have approved the manuscript and agree with its submission. The authors declare no
competing interests.

Thank you for considering our submission.

Sincerely,

Almaz Sapargali (corresponding author)\\
almaz.sapargali2001@gmail.com

\end{document}
